%% ============================================================================
\section{Lecture 7: Convection, Atmospheric Stability, and Dimensionless Control}
\label{sec:lec07}
%% ============================================================================
\date{28/05/2026}

\begin{center}
    \textit{Lecture date: 28/05/2026}
\end{center}

The previous lecture built a unified picture of transport by microscopic
random walks: diffusion of particles, diffusion of heat, and diffusion of
momentum. This lecture turns to a different kind of heat transport:
\emph{convection}. Instead of heat being carried by microscopic steps, a whole
fluid element moves, expands or contracts, and mixes with its surroundings.
Gravity is the essential new ingredient, because it lets density differences
turn into buoyant acceleration.

%%----------------------------------------------------------------------------
\subsection{Three Ways to Transport Heat}
\label{sec:lec07:heat-transport-modes}
%%----------------------------------------------------------------------------

\subsubsection*{Physical Motivation}
Before asking when convection starts, it is useful to separate it from the
other two heat-transfer mechanisms. A warm object in a room can exchange heat
with its surroundings by:
\begin{itemize}
    \item \textbf{conduction}: molecular collisions transfer energy through
    matter, described by Fourier's law;
    \item \textbf{radiation}: photons carry energy, even through vacuum;
    \item \textbf{convection}: the material itself moves, so hot and cold
    parcels are physically exchanged.
\end{itemize}
Convection is therefore not just ``better conduction''. It changes the
configuration of the fluid.

In the Earth's atmosphere, sunlight mostly reaches the surface directly because
the atmosphere is relatively transparent at visible wavelengths. The warmed
surface then emits infrared radiation, to which the atmosphere is much more
opaque. The lower atmosphere is heated from below, and the question becomes:
will the heat escape by slow conduction/radiation, or will the air itself begin
to rise and overturn?

\dfn[def:lec07:convection]{Convection}{
Convection is heat transport by the bulk motion of a fluid. In a gravitational
field it is driven by buoyancy: a parcel whose density is lower than the
density of its surroundings accelerates upward, while a denser parcel
accelerates downward.
}

%%----------------------------------------------------------------------------
\subsection{Hydrostatic Balance and an Isothermal Atmosphere}
\label{sec:lec07:hydrostatic-isothermal}
%%----------------------------------------------------------------------------

\subsubsection*{Mathematical Setup}
Let $z$ denote height, $P(z)$ pressure, $\rho(z)$ mass density, $T(z)$
temperature, $g$ the gravitational acceleration, and $\mu$ the mass per gas
particle. A static atmosphere must balance gravity against the pressure
gradient. Per unit mass,
\begin{equation}
    \boxed{
    \frac{1}{\rho}\frac{dP}{dz} = -g .
    }
    \label{eq:lec07:hydrostatic-balance}
\end{equation}
\textit{Significance:} Pressure must decrease upward because lower layers
support the weight of all the gas above them.

For an ideal gas,
\begin{equation}
    P = n k_B T = \frac{\rho}{\mu}k_B T ,
    \label{eq:lec07:ideal-gas}
\end{equation}
where $n=\rho/\mu$ is the particle number density. Equations
\eqref{eq:lec07:hydrostatic-balance} and \eqref{eq:lec07:ideal-gas} are not
yet enough to determine the atmosphere, because both $\rho(z)$ and $T(z)$ are
unknown. One possible closure is an isothermal atmosphere, $T(z)=T_0$.

Substituting $P=\rho k_B T_0/\mu$ into hydrostatic balance gives
\begin{align}
    \frac{1}{\rho}\frac{d}{dz}
    \left(\frac{\rho k_B T_0}{\mu}\right)
    &= -g ,
    \label{eq:lec07:isothermal-step1}\\
    \frac{k_B T_0}{\mu}\frac{1}{\rho}\frac{d\rho}{dz}
    &= -g ,
    \label{eq:lec07:isothermal-step2}\\
    \frac{d\ln\rho}{dz}
    &= -\frac{\mu g}{k_B T_0}.
    \label{eq:lec07:isothermal-step3}
\end{align}
Integrating from $z=0$ to $z$,
\begin{equation}
    \boxed{
    \rho(z)=\rho_0 e^{-z/H},
    \qquad
    H=\frac{k_B T_0}{\mu g}.
    }
    \label{eq:lec07:isothermal-atmosphere}
\end{equation}
\textit{Significance:} $H$ is the atmospheric scale height: it is the height
over which the density falls by a factor of $e$.

For Earth, taking $T_0\simeq300\,\si{K}$, $\mu\simeq 29m_p$, and
$g\simeq10\,\si{m\,s^{-2}}$ gives
\begin{equation}
    H \sim 8\,\si{km}.
    \label{eq:lec07:earth-scale-height}
\end{equation}
This does not mean that the atmosphere ends after $8\,\si{km}$; it means that
the density has dropped by one $e$-fold.

%%----------------------------------------------------------------------------
\subsection{Adiabatic Atmospheres and the Lapse Rate}
\label{sec:lec07:adiabatic-atmosphere}
%%----------------------------------------------------------------------------

\subsubsection*{Conceptual Background}
Convection is tested by displacing a fluid element. Pressure equilibrates
quickly, on roughly a sound-crossing time, while heat exchange by conduction or
radiation can be much slower. Therefore a rapidly displaced parcel is often
approximately \emph{adiabatic}: it changes pressure, density, and temperature
without exchanging heat with its surroundings.

For an ideal gas undergoing adiabatic changes,
\begin{equation}
    P = K\rho^\gamma ,
    \label{eq:lec07:adiabatic-relation}
\end{equation}
where $K$ is constant along the parcel trajectory and
$\gamma=c_p/c_V$ is the adiabatic index. Combining this with
$P=\rho k_B T/\mu$ gives
\begin{equation}
    T \propto \rho^{\gamma-1},
    \qquad
    P \propto T^{\gamma/(\gamma-1)} .
    \label{eq:lec07:T-rho-adiabatic}
\end{equation}

Differentiate the ideal-gas pressure relation along an adiabatic profile:
\begin{align}
    dP &= \gamma\frac{P}{\rho}\,d\rho ,
    \label{eq:lec07:adiabatic-dP}\\
    dT &= \frac{\mu}{k_B}\left(\frac{dP}{\rho}
    - \frac{P}{\rho^2}\,d\rho\right).
    \label{eq:lec07:dT-general}
\end{align}
Using $d\rho = (\rho/\gamma P)dP$ in Eq.~\eqref{eq:lec07:dT-general},
\begin{align}
    dT
    &= \frac{\mu}{k_B}\left(\frac{1}{\rho}
    - \frac{1}{\gamma\rho}\right)dP
    = \frac{\mu}{k_B}\frac{\gamma-1}{\gamma}\frac{dP}{\rho}.
    \label{eq:lec07:dT-adiabatic}
\end{align}
Hydrostatic balance gives $dP/dz=-\rho g$, hence
\begin{equation}
    \boxed{
    \left(\frac{dT}{dz}\right)_{\rm ad}
    = -\,\frac{\gamma-1}{\gamma}\,\frac{\mu g}{k_B}
    = -\,\frac{g}{c_p}.
    }
    \label{eq:lec07:dry-adiabatic-lapse}
\end{equation}
\textit{Significance:} A dry parcel cools as it rises because it expands and
does work on its surroundings. The atmospheric temperature gradient cannot be
arbitrarily steep without triggering convection.

For dry air, this gives roughly
\begin{equation}
    \left|\frac{dT}{dz}\right|_{\rm ad}
    \sim 10\,\si{K\,km^{-1}}
    \simeq 1\,^\circ\mathrm{C}/100\,\si{m}.
    \label{eq:lec07:dry-lapse-number}
\end{equation}

\nt{Moist air has a smaller effective adiabatic lapse rate. As a rising parcel
expands and cools, water vapor can condense and release latent heat. That heat
partly offsets the cooling, so saturated air can become convectively unstable
for a weaker temperature gradient than dry air.}

%%----------------------------------------------------------------------------
\subsection{Convective Stability as an Entropy Criterion}
\label{sec:lec07:convective-stability}
%%----------------------------------------------------------------------------

\subsubsection*{Physical Motivation}
It is tempting to say that an atmosphere is stable because dense air is below
light air. That is not sufficient. A displaced parcel does not keep its
original density: it expands or contracts to match the ambient pressure. The
right comparison is made after the parcel has adjusted its pressure but before
it has exchanged heat. In other words, the parcel keeps its entropy.

For an ideal gas, the entropy per unit mass can be written, up to an additive
constant, as
\begin{equation}
    s = c_V\ln\!\left(\frac{P}{\rho^\gamma}\right)
    = c_p\ln T - \frac{k_B}{\mu}\ln P + \text{const}.
    \label{eq:lec07:entropy-ideal-gas}
\end{equation}
An adiabatic displacement keeps $s$ fixed. Suppose a parcel from height $z$ is
lifted slightly to $z+dz$. It rapidly adjusts to the new pressure
$P(z+dz)$, but retains $s(z)$. Then:
\begin{itemize}
    \item if the environmental entropy increases upward, $ds/dz>0$, the
    lifted parcel has lower entropy than its new surroundings; at the same
    pressure it is denser, so it sinks back;
    \item if the environmental entropy decreases upward, $ds/dz<0$, the
    lifted parcel has higher entropy than its new surroundings; at the same
    pressure it is lighter, so it keeps rising.
\end{itemize}
Thus
\begin{equation}
    \boxed{
    \frac{ds}{dz}>0 \quad \Longrightarrow \quad \text{stable},
    \qquad
    \frac{ds}{dz}<0 \quad \Longrightarrow \quad \text{convectively unstable}.
    }
    \label{eq:lec07:entropy-stability}
\end{equation}
\textit{Significance:} Convection begins when the one-dimensional stratified
solution is unstable to vertical displacements; the system then develops
overturning motions instead of staying purely radial or vertical.

Using Eq.~\eqref{eq:lec07:dry-adiabatic-lapse}, the same criterion can be
phrased as a temperature-gradient condition. Define the environmental lapse
rate
\begin{equation}
    \Gamma \equiv -\frac{dT}{dz},
    \qquad
    \Gamma_{\rm ad} \equiv \frac{g}{c_p}.
    \label{eq:lec07:lapse-definitions}
\end{equation}
Then
\begin{equation}
    \boxed{
    \Gamma < \Gamma_{\rm ad} \quad \Longrightarrow \quad \text{stable},
    \qquad
    \Gamma > \Gamma_{\rm ad} \quad \Longrightarrow \quad \text{convectively unstable}.
    }
    \label{eq:lec07:lapse-stability}
\end{equation}
\textit{Significance:} A temperature that falls with height faster than the
adiabatic rate makes rising parcels remain warmer and lighter than their
surroundings.

In ordinary dry air, the condition says that height alone should not produce a
temperature difference much larger than about $1^\circ\mathrm{C}$ per
$100\,\si{m}$. If the gradient becomes steeper, convection mixes the air and
drives the profile back toward marginal stability.

%%----------------------------------------------------------------------------
\subsection{Why Efficient Convection Drives Marginal Stability}
\label{sec:lec07:marginal-stability}
%%----------------------------------------------------------------------------

\subsubsection*{Physical Motivation}
Convection often looks like a violent instability, but in many systems it
self-regulates. If the entropy gradient is strongly unstable, buoyant motions
become fast and carry heat very efficiently. That heat transport reduces the
entropy contrast that caused the motion. The final state is usually close to
the boundary between stable and unstable.

This explains why convective regions in stars, atmospheres, and planetary
interiors are often nearly adiabatic. They are not exactly adiabatic, because
a small superadiabatic excess is needed to carry the heat flux, but the excess
can be tiny when convection is efficient.

Examples discussed in the lecture:
\begin{itemize}
    \item \textbf{Earth's atmosphere:} surface heating can make the lower
    atmosphere unstable; overturning then limits the lapse rate.
    \item \textbf{Earth's mantle:} the mantle can convect and carry heat from
    the hot interior, while the rigid crust conducts heat across the final
    boundary layer.
    \item \textbf{the Sun:} convection in the outer envelope transports heat
    by bulk motion; if the whole Sun were strongly superadiabatic, it would
    mix on a sound-crossing timescale of only hours, so the actual convective
    stratification is close to marginal.
\end{itemize}

%%----------------------------------------------------------------------------
\subsection{Buoyancy and a Convective Velocity Scale}
\label{sec:lec07:buoyancy-velocity}
%%----------------------------------------------------------------------------

\subsubsection*{Mathematical Setup}
Let $\alpha$ be the thermal expansion coefficient,
\begin{equation}
    \alpha \equiv -\frac{1}{\rho}
    \left(\frac{\partial \rho}{\partial T}\right)_P .
    \label{eq:lec07:thermal-expansion}
\end{equation}
For a small temperature excess $\Delta T$ at fixed pressure,
\begin{equation}
    \frac{\Delta\rho}{\rho}\simeq -\alpha\,\Delta T .
    \label{eq:lec07:density-temperature}
\end{equation}
For an ideal gas, $\alpha=1/T$.

Consider a parcel with density $\rho_1$ embedded in surroundings of density
$\rho_2$. The gravitational force on a volume $V$ is $\rho_1 Vg$ downward,
while the buoyant force is $\rho_2 Vg$ upward. The net acceleration is
\begin{equation}
    a
    = g\,\frac{\rho_2-\rho_1}{\rho_1}
    \sim g\,\frac{\Delta\rho}{\rho}.
    \label{eq:lec07:buoyant-acceleration-density}
\end{equation}
Using Eq.~\eqref{eq:lec07:density-temperature},
\begin{equation}
    \boxed{
    a_{\rm buoy} \sim g\,\alpha\,\Delta T_{\rm ex}.
    }
    \label{eq:lec07:buoyant-acceleration}
\end{equation}
\textit{Significance:} Only the temperature excess beyond the adiabatic
background contributes to buoyant acceleration; the ordinary adiabatic cooling
of a lifted parcel is already part of the neutral reference state.

If the unstable layer has height $H$, the buoyancy time and the corresponding
free convective velocity scale are
\begin{align}
    \boxed{
    \tau_{\rm buoy}
    \sim \sqrt{\frac{H}{g\alpha\Delta T_{\rm ex}}},
    }
    \label{eq:lec07:buoyancy-time}\\
    \boxed{
    U_{\rm conv}
    \sim \frac{H}{\tau_{\rm buoy}}
    \sim \sqrt{g\alpha\Delta T_{\rm ex}H}.
    }
    \label{eq:lec07:convective-velocity}
\end{align}
\textit{Significance:} Larger gravity, stronger thermal expansion, larger
superadiabatic temperature contrast, and taller layers all make convection
faster.

%%----------------------------------------------------------------------------
\subsection{Dimensionless Numbers for Convection}
\label{sec:lec07:dimensionless-convection}
%%----------------------------------------------------------------------------

\subsubsection*{Motivation}
The final part of the lecture returns to dimensional analysis. A convective
layer is controlled by the layer height $H$, the superadiabatic temperature
contrast $\Delta T_{\rm ex}$, gravity $g$, thermal expansion $\alpha$,
kinematic viscosity $\nu$, and thermal diffusivity $\chi$. Both $\nu$ and
$\chi$ have units of $\si{m^2\,s^{-1}}$, but they diffuse different quantities:
momentum and heat.

\dfn[def:lec07:thermal-diffusivity]{Thermal Diffusivity}{
The thermal diffusivity $\chi$ is the diffusion coefficient for temperature:
\[
    \frac{\partial T}{\partial t}=\chi\nabla^2T .
\]
It should be distinguished from particle diffusion and from kinematic
viscosity, even though all three have dimensions of length squared over time.
}

The first dimensionless number is the Prandtl number:
\begin{equation}
    \boxed{
    \mathrm{Pr}=\frac{\nu}{\chi}.
    }
    \label{eq:lec07:prandtl}
\end{equation}
\textit{Significance:} $\mathrm{Pr}$ compares how fast momentum diffuses to
how fast heat diffuses.

The second important comparison is between buoyant driving and diffusive
damping. The heat-diffusion time across the layer is
\begin{equation}
    \tau_\chi \sim \frac{H^2}{\chi},
    \label{eq:lec07:thermal-diffusion-time}
\end{equation}
and the viscous momentum-diffusion time is
\begin{equation}
    \tau_\nu \sim \frac{H^2}{\nu}.
    \label{eq:lec07:viscous-diffusion-time}
\end{equation}
Combining these with the buoyancy time
\eqref{eq:lec07:buoyancy-time} gives the Rayleigh number:
\begin{equation}
    \boxed{
    \mathrm{Ra}
    = \frac{g\alpha\Delta T_{\rm ex}H^3}{\nu\chi}.
    }
    \label{eq:lec07:rayleigh}
\end{equation}
\textit{Significance:} $\mathrm{Ra}$ measures buoyant driving against the two
diffusive processes that suppress convection: viscosity damps velocity, and
thermal diffusion erases the temperature anomaly that produces buoyancy.

Equivalently,
\begin{equation}
    \mathrm{Ra}
    = \frac{\tau_\nu\,\tau_\chi}{\tau_{\rm buoy}^2}.
    \label{eq:lec07:rayleigh-timescales}
\end{equation}
Large $\mathrm{Ra}$ means buoyancy acts before diffusion can smooth the
temperature and momentum perturbations. Small $\mathrm{Ra}$ means diffusion
wins, and the layer remains conductive. In the standard Rayleigh-Benard
problem with fixed-temperature plates, convection starts at
$\mathrm{Ra}\simeq1708$; for order-of-magnitude modeling the important point is
the scaling $\mathrm{Ra}\gtrsim 1$ versus $\mathrm{Ra}\ll1$.

\subsubsection*{Relation to Reynolds number}
Using the convective velocity estimate
\eqref{eq:lec07:convective-velocity}, the Reynolds number of the convective
motion is
\begin{equation}
    \mathrm{Re}_{\rm conv}
    = \frac{U_{\rm conv}H}{\nu}
    \sim \frac{H}{\nu}\sqrt{g\alpha\Delta T_{\rm ex}H}.
    \label{eq:lec07:convective-reynolds}
\end{equation}
Squaring and using Eq.~\eqref{eq:lec07:rayleigh},
\begin{equation}
    \boxed{
    \mathrm{Re}_{\rm conv}^2
    \sim \frac{\mathrm{Ra}}{\mathrm{Pr}}.
    }
    \label{eq:lec07:re-ra-pr}
\end{equation}
\textit{Significance:} The speed of convective stirring is not an externally
given input like the speed of a stirred cup; it is part of the solution, set by
buoyancy and limited by diffusion.

%%----------------------------------------------------------------------------
\subsection{Summary of the Lecture}
\label{sec:lec07:summary}
%%----------------------------------------------------------------------------

\begin{itemize}
    \item Hydrostatic balance fixes $dP/dz=-\rho g$, but does not by itself
    determine $T(z)$ and $\rho(z)$.
    \item An isothermal ideal-gas atmosphere has
    $\rho=\rho_0e^{-z/H}$ with $H=k_BT/(\mu g)$.
    \item A rapidly displaced parcel evolves adiabatically, so the neutral
    temperature gradient is $dT/dz=-g/c_p$.
    \item The stability criterion is entropy stratification:
    $ds/dz>0$ is stable, while $ds/dz<0$ is convectively unstable.
    \item The practical atmospheric form is
    $\Gamma<\Gamma_{\rm ad}$ stable and $\Gamma>\Gamma_{\rm ad}$ unstable,
    with $\Gamma_{\rm ad}\simeq10\,\si{K\,km^{-1}}$ for dry air.
    \item Buoyant acceleration is
    $a_{\rm buoy}\sim g\alpha\Delta T_{\rm ex}$, where
    $\Delta T_{\rm ex}$ is the superadiabatic temperature excess.
    \item The main dimensionless control parameters are
    $\mathrm{Pr}=\nu/\chi$ and
    $\mathrm{Ra}=g\alpha\Delta T_{\rm ex}H^3/(\nu\chi)$.
\end{itemize}
